\section{8.3 중간고사 대비 총정리}
\begin{frame}{이 절: 새 개념 0개, 하는 일 세 가지}
  \begin{itemize}
    \item (1) Ch\,2$\sim$7 이론을 \textbf{``손으로 재현할 수 있는
      형태''}로 압축.
    \item (2) 준비도를 스스로 재는 \textbf{자가진단} 제시.
    \item (3) 8.1에서 4주간 쌓아온 프로젝트 절차를 복습 자료로 재작성.
  \end{itemize}
  \vskip0.6em
  \begin{block}{시험의 성격}
    시험은 \textbf{``Block A의 개념을 \emph{새로운 MDP}에 적용하는
    것''}이므로, \textbf{암기는 쓸모가 없고} 이 절은
    \textbf{``스스로 유도하는 습관''} 하나에 맞추어져 있다.
  \end{block}
  \vskip0.6em
  {\footnotesize 특히 ``벨만방정식을 터미널부터 거꾸로 대입해 직접
  풀기''(Ch\,4)와 ``TD 오차의 수렴 조건 논증''(Ch\,6)이 다른 MDP에도
  그대로 적용되는지가 좋은 복습.}
\end{frame}
